% !TEX root = ../main.tex
%% Sujet type bac n° 3 — sujet inédit OptiTAB (calculs vérifiés par programme : bac_s3.py)

\sujetBac{3}{Probabilités et concentration/4,Équation différentielle/5,Vrai ou faux/4,Fonction logarithme et intégration par parties/5}

\begin{exoBac}{4}{Probabilités et concentration}
Une entreprise achète ses batteries à deux fournisseurs : $60\,\%$ au fournisseur A et $40\,\%$ au fournisseur B. On sait que $3\,\%$
des batteries du fournisseur A sont défectueuses, contre $8\,\%$ de celles du fournisseur B. On prélève une batterie au hasard dans le
stock et l'on note $A$ l'événement « la batterie provient du fournisseur A » et $D$ l'événement « la batterie est défectueuse ».
\partieExo{A}{Une batterie au hasard}
\begin{questions}
  \item Représenter la situation par un arbre pondéré.
  \item Montrer que $P(D)=0{,}05$.
  \item La batterie prélevée est défectueuse. Quelle est la probabilité qu'elle provienne du fournisseur B ?
  \item Les événements $A$ et $D$ sont-ils indépendants ? Justifier.
\end{questions}
\partieExo{B}{Un lot de 200 batteries}
On prélève un lot de $200$ batteries ; le stock est assez important pour assimiler ce prélèvement à $200$ tirages indépendants avec
remise. On note $X$ le nombre de batteries défectueuses du lot. On admet que $X$ suit la loi binomiale $\mathcal{B}(200\pv0{,}05)$.
\begin{questions}[resume]
  \item Calculer l'espérance et la variance de $X$.
  \item Calculer la probabilité que le lot contienne au plus $12$ batteries défectueuses. On arrondira au millième.
  \item À l'aide de l'inégalité de Bienaymé-Tchebychev, montrer que $P\big(5<X<15\big)\geq0{,}62$.
  \item Calculer la valeur exacte arrondie au millième de $P\big(5<X<15\big)$ et commenter.
\end{questions}
\end{exoBac}

\begin{exoBac}{5}{Équation différentielle}
Un café est servi à $85\,^\circ$C dans une pièce dont la température est constante, égale à $20\,^\circ$C. On note $\theta(t)$ la
température du café, en degrés Celsius, $t$ minutes après le service. On admet que la fonction $\theta$ est dérivable sur
$\intFO0{+\infty}$, que $\theta(0)=85$ et qu'elle vérifie, selon la loi de refroidissement de Newton :
\[(E):\qquad \theta'(t)=-0{,}08\big(\theta(t)-20\big).\]
\begin{questions}
  \item Vérifier que l'équation $(E)$ s'écrit $y'=-0{,}08\,y+1{,}6$, puis déterminer la solution constante de $(E)$.
  \item Résoudre l'équation différentielle $(E)$, puis montrer que, pour tout $t\geq0$, $\theta(t)=20+65\,e^{-0{,}08t}$.
  \item Étudier le sens de variation de $\theta$ et déterminer sa limite en $+\infty$. Interpréter.
  \item Le café est jugé buvable lorsque sa température descend à $60\,^\circ$C. Au bout de combien de temps, à la seconde près,
    est-il buvable ?
  \item Au bout de combien de temps l'écart entre la température du café et celle de la pièce est-il divisé par deux ?
  \item Calculer la température moyenne du café pendant les dix premières minutes, c'est-à-dire
    $\dfrac{1}{10}\displaystyle\int_0^{10}\theta(t)\,\mathrm{d}t$. On arrondira au dixième de degré.
\end{questions}
\end{exoBac}

\begin{exoBac}{4}{Vrai ou faux}
Pour chacune des affirmations suivantes, indiquer si elle est vraie ou fausse, en justifiant la réponse. Une réponse non justifiée
n'est pas prise en compte. Dans les affirmations 1 et 2, l'espace est muni d'un repère orthonormé.
\begin{questions}
  \item \textbf{Affirmation 1.} La droite $d$ de représentation paramétrique $\begin{cases}x=1+2t\\y=-t\\z=3+t\end{cases}$, $t\in\R$, est
    strictement parallèle au plan $\mathcal{P}$ d'équation $x+y-z+4=0$.
  \item \textbf{Affirmation 2.} Les droites $d_1:\begin{cases}x=t\\y=1+t\\z=2t\end{cases}$, $t\in\R$, et
    $d_2:\begin{cases}x=1+s\\y=3-s\\z=s\end{cases}$, $s\in\R$, sont sécantes.
  \item \textbf{Affirmation 3.} Dans une association de $10$ membres, il y a $120$ façons d'élire un bureau composé d'un président,
    d'un trésorier et d'un secrétaire, trois personnes distinctes.
  \item \textbf{Affirmation 4.} La suite $(u_n)$ définie pour tout entier naturel $n$ par $u_n=\dfrac{2n+1}{n+3}$ est croissante et
    converge vers $2$.
\end{questions}
\end{exoBac}

\begin{exoBac}{5}{Fonction logarithme et intégration par parties}
On considère la fonction $f$ définie sur $\intOO0{+\infty}$ par $f(x)=\dfrac{\ln x}{x^{2}}$, et l'on note $\mathcal{C}_f$ sa courbe
représentative dans un repère orthonormé.
\partieExo{A}{Étude de la fonction}
\begin{questions}
  \item Déterminer les limites de $f$ en $0$ et en $+\infty$. En déduire les asymptotes éventuelles à $\mathcal{C}_f$.
  \item Montrer que, pour tout réel $x>0$, $f'(x)=\dfrac{1-2\ln x}{x^{3}}$. Dresser le tableau de variations de $f$, en précisant la
    valeur exacte de son maximum.
  \item Résoudre l'équation $f(x)=0$, puis étudier le signe de $f$ sur $\intOO0{+\infty}$.
\end{questions}
\partieExo{B}{Calculs d'aires}
\begin{questions}[resume]
  \item À l'aide d'une intégration par parties, calculer $I=\displaystyle\int_1^{e}\frac{\ln x}{x^{2}}\,\mathrm{d}x$.
  \item Interpréter graphiquement le nombre $I$.
  \item Pour tout réel $\lambda>1$, on pose $A(\lambda)=\displaystyle\int_1^{\lambda}f(x)\,\mathrm{d}x$. Montrer que
    $A(\lambda)=1-\dfrac{1+\ln\lambda}{\lambda}$, puis déterminer la limite de $A(\lambda)$ quand $\lambda$ tend vers $+\infty$.
\end{questions}
\end{exoBac}
